\chapter{Graphs through Their Eigenvalues}\label{linear:ch:graphs}

\begin{center}\itshape
In this paper, we show that every $(2^{n-1}+1)$-vertex induced subgraph of the $n$-dimensional hypercube has maximum degree at
least $\sqrt n$.
\upshape(Hao Huang, opening of the abstract of his 2019 paper in the \emph{Annals of Mathematics})
\end{center}

Here is a puzzle that Paul Erd\H{o}s, Alfr\'ed R\'enyi and Vera T.~S\'os solved in 1966. At a party, every two guests have
exactly one common friend among the guests. Prove that some guest is a friend of everybody. (Friendship is mutual, and
nobody is their own friend.) The answer, the \emph{friendship theorem}, is a little surprising, and its proof is even more so:
the decisive step is that a certain number would have to be the square root of an integer and rational at the same time.
The party is a graph, and the argument runs through the eigenvalues of its adjacency matrix. Martin Aigner and G\"unter
Ziegler put it in \emph{Proofs from THE BOOK} under the title ``Of friends and politicians'': the guest who knows everybody
is the politician.

Fifty-three years later, on 1 July 2019, Hao Huang of Emory University posted a six-page preprint proving the
\emph{sensitivity conjecture}, a question about Boolean functions that had resisted computer scientists since 1992 and
that many experts expected to need heavy machinery. The proof fits on a page. It builds a $2^n\times2^n$ matrix with entries
$0,\pm1$ whose square is $n$ times the identity, and applies Cauchy's interlacing theorem from Chapter~\ref{linear:ch:spectral}.
The next day, the computer scientist Scott Aaronson, who had called the conjecture ``one of the most frustrating and
embarrassing open problems in all of combinatorics and theoretical computer science'', wrote that reading and understanding
the ``elementary 1.5-page argument'' had taken him about half an hour.

This chapter is about such arguments. The adjacency matrix of a graph is symmetric, so it has real eigenvalues and an
orthonormal eigenbasis, and every theorem of Chapters~\ref{linear:ch:spectral} and~\ref{linear:ch:perron} becomes a statement about graphs.
We compute the spectra of the basic families, prove the friendship theorem, classify the possible Moore graphs (a story with
an unsolved case), prove Hoffman's bound for independent sets, the Graham--Pollak theorem, and Huang's theorem. Other parts
of spectral graph theory live elsewhere in the series: the matrix-tree theorem, Cheeger's inequality and expanders in
Tome~11, Appendix~A; Ramanujan graphs in Tome~11.

\begin{think}
Try the party puzzle on small cases before reading on. With three guests, a triangle works. Can you find a party with five
guests? With seven? Do all your examples have a guest who knows everybody? Is there any example with no such guest and with
everybody having the same number of friends?
\end{think}

\section{Adjacency matrices and walks}

A graph $G$ has a finite vertex set $V$, $|V|=n$, and a set of edges, unordered pairs of distinct vertices; we write $i\sim j$ for
an edge. Its \emph{adjacency matrix} $A$ has $a_{ij}=1$ if $i\sim j$ and $0$ otherwise. It is real symmetric with zero diagonal,
so its eigenvalues are real, $\lambda_1\ge\dots\ge\lambda_n$, and their sum is $\operatorname{tr}A=0$. The spectrum does not
depend on the numbering of the vertices, since renumbering conjugates $A$ by a permutation matrix.

As in Chapter~\ref{linear:ch:perron}, $(A^k)_{ij}$ counts the walks of length $k$ from $i$ to $j$. Hence $\operatorname{tr}A^2=\sum_i\deg i
=2|E|$ and $\operatorname{tr}A^3=6\,t(G)$, where $t(G)$ is the number of triangles (each triangle gives six closed walks of length
$3$). The spectrum therefore knows the number of vertices, edges and triangles.

If $G$ is \emph{$k$-regular} (all degrees equal $k$), then $A\mathbf 1=k\mathbf 1$, and by Exercise~\ref{linear:ex:rowsum}, $k=\lambda_1$
is the largest eigenvalue; Perron--Frobenius (Theorem~\ref{linear:thm:frobenius}) says that $G$ is connected if and only if $k$ is a
simple eigenvalue (a disconnected graph has an eigenvector $\mathbf 1_C$ for each component $C$). For regular graphs, the other
eigenvectors can be taken orthogonal to $\mathbf 1$, on which the all-ones matrix $J$ acts as $0$: this is why regular graphs are
so much easier.

\begin{example}[basic spectra]\label{linear:exa:spectra}
\begin{enumerate}[label=(\alph*)]
\item The complete graph $K_n$ has $A=J-I$, with eigenvalues $n-1$ (once) and $-1$ ($n-1$ times).
\item The cycle $C_n$ has $A=S+S^{-1}$ with $S$ the cyclic shift, so its eigenvalues are $\omega^j+\omega^{-j}=2\cos\frac{2\pi j}n$,
$j=0,\dots,n-1$, with eigenvectors $(\omega^{jm})_m$, $\omega=e^{2\pi i/n}$.
\item The path $P_n$ has eigenvalues $2\cos\frac{j\pi}{n+1}$, $j=1,\dots,n$ (Exercise~\ref{linear:ex:pathgraph}).
\item The hypercube $Q_n$, with vertex set $\{0,1\}^n$ and edges between words differing in one coordinate, has eigenvalues
$n-2j$ with multiplicity $\binom nj$, $j=0,\dots,n$ (Exercise~\ref{linear:ex:hypercube}).
\end{enumerate}
\end{example}

\section{The friendship theorem}

\begin{theorem}[Erd\H{o}s--R\'enyi--S\'os, 1966]\label{linear:thm:friendship}
Let $G$ be a finite graph in which every two distinct vertices have exactly one common neighbour. Then $G$ is a
\emph{windmill}: $n=2m+1$ and $G$ consists of $m$ triangles sharing one common vertex.
\end{theorem}

\begin{figure}[ht]
\centering
\fig{linear-ch08-windmill}
\caption{The windmill with four blades: every two vertices have exactly one common neighbour, and the centre is the
politician.}
\label{linear:fig:windmill}
\end{figure}

\begin{proof}
\emph{Step 1: non-adjacent vertices have the same degree.} Let $u\not\sim v$. For each neighbour $x$ of $u$, let $f(x)$ be the
unique common neighbour of $x$ and $v$; it is a neighbour of $v$. If $f(x)=f(x')=y$ for neighbours $x\ne x'$ of $u$, then $u$ and
$y$ have the two common neighbours $x$ and $x'$ (and $y\ne u$, since $y\sim v$ and $u\not\sim v$), which is forbidden. So $f$ is
injective and $\deg u\le\deg v$; by symmetry, $\deg u=\deg v$.

\emph{Step 2: if no vertex is adjacent to all others, $G$ is regular.} By Step~1, the degree is constant on each connected component
of the complementary graph $\bar G$ (where $u,v$ are joined when $u\not\sim v$ in $G$). If $\bar G$ is disconnected, split $V=V_1\sqcup
V_2$ with no edge of $\bar G$ between $V_1$ and $V_2$: every vertex of $V_1$ is adjacent in $G$ to every vertex of $V_2$. If both parts
had at least two vertices, two vertices of $V_1$ would have two common neighbours in $V_2$. So one part is a single vertex, adjacent
to all others. Hence, if there is no such vertex, $\bar G$ is connected and $G$ is regular.

\emph{Step 3: the regular case is impossible.} Let $G$ be $k$-regular. The hypothesis says $(A^2)_{uv}=1$ for $u\ne v$, and
$(A^2)_{uu}=k$, so
\[
A^2=(k-1)I+J.
\]
Counting walks of length $2$ from a vertex, $k(k-1)=n-1$ (every other vertex is reached exactly once, and each walk
$u\to x\to y$ with $y\ne u$ reaches a vertex other than $u$). The matrix $A^2$ has eigenvalue $k-1+n=k^2$ on $\mathbf 1$ and $k-1$ on
$\mathbf 1^\perp$. So $A$ has the eigenvalue $k$ on $\mathbf 1$ and eigenvalues $\pm\sqrt{k-1}$ on $\mathbf 1^\perp$, say $+\sqrt{k-1}$
with multiplicity $r$ and $-\sqrt{k-1}$ with multiplicity $s$. The trace is zero:
\[
k+(r-s)\sqrt{k-1}=0.
\]
Here $k\ge1$, since two vertices need a common neighbour, so the equation forces $r\ne s$. Then
$\sqrt{k-1}=\frac k{s-r}$ is rational, hence an integer $h$ (a rational square root of an integer is an integer). Then $h(s-r)=h^2+1$,
so $h$ divides $1$, $h=1$, $k=2$ and $n=k(k-1)+1=3$: $G$ is a triangle. But in a triangle every vertex is adjacent to all others,
contradicting the assumption of Step 3.

\emph{Step 4: the structure.} So some vertex $w$ is adjacent to all others. For any other vertex $u$, the common neighbour of $u$
and $w$ is a vertex $u'\ne w$ adjacent to $u$; and $u$ has no other neighbour besides $w$ and $u'$, for a second one $u''$ would be
a second common neighbour of $u$ and $w$. So $V\setminus\{w\}$ is partitioned into pairs $\{u,u'\}$, each forming a triangle with $w$.
\end{proof}

Compare the proof with the statement: the conclusion is a structure, but the obstacle in Step~3 is arithmetic, the irrationality
of $\sqrt{k-1}$. This pattern, ``the eigenvalues would have to be integers, and they are not'', returns twice below. There is no
known proof of the friendship theorem without some version of this counting; it famously fails for infinite graphs, where
friendship graphs without a politician exist (Exercise~\ref{linear:ex:infinitefriends}).

\section{Strongly regular graphs and Moore graphs}

\begin{definition}\label{linear:def:srg}
A $k$-regular graph on $n$ vertices is \emph{strongly regular} with parameters $(n,k,\lambda,\mu)$ if every two adjacent
vertices have exactly $\lambda$ common neighbours and every two non-adjacent distinct vertices have exactly $\mu$.
\end{definition}

The friendship hypothesis would be $\lambda=\mu=1$. The \emph{Petersen graph} (Figure~\ref{linear:fig:petersen}) is strongly regular with
parameters $(10,3,0,1)$; so is the pentagon, $(5,2,0,1)$. In matrix form the definition says
\begin{equation}\label{linear:eq:srg}
A^2=kI+\lambda A+\mu(J-I-A).
\end{equation}

\begin{figure}[ht]
\centering
\fig{linear-ch08-petersen}
\caption{The Petersen graph, strongly regular with parameters $(10,3,0,1)$: no triangles, and two non-adjacent vertices have
exactly one common neighbour. Its spectrum is $3$, $1$ (five times), $-2$ (four times). Its vertices can also be read as the
$2$-element subsets of $\{1,\dots,5\}$, adjacent when disjoint.}
\label{linear:fig:petersen}
\end{figure}

\begin{theorem}[eigenvalues of strongly regular graphs]\label{linear:thm:srg}
Let $G$ be connected and strongly regular with parameters $(n,k,\lambda,\mu)$, $\mu>0$. Its eigenvalues are $k$ (simple) and the
two roots $\theta>\tau$ of $x^2-(\lambda-\mu)x-(k-\mu)=0$, with multiplicities
\[
m_{\theta},m_\tau=\frac12\Bigl((n-1)\mp\frac{2k+(n-1)(\lambda-\mu)}{\sqrt{(\lambda-\mu)^2+4(k-\mu)}}\Bigr).
\]
In particular these two numbers must be non-negative integers.
\end{theorem}
\begin{proof}
On $\mathbf 1^\perp$, which is invariant and spanned by eigenvectors, $J=0$ and~\eqref{linear:eq:srg} reads $A^2=(k-\mu)I+(\lambda-\mu)A$. So
every eigenvalue $x$ of $A$ on $\mathbf 1^\perp$ satisfies $x^2-(\lambda-\mu)x-(k-\mu)=0$. The multiplicities satisfy $m_\theta+m_\tau=n-1$
and, since the trace is $0$, $k+m_\theta\theta+m_\tau\tau=0$. With $\theta+\tau=\lambda-\mu$ and $\theta-\tau=\sqrt\Delta$,
$\Delta=(\lambda-\mu)^2+4(k-\mu)$, the second equation becomes $k+(n-1)\frac{\lambda-\mu}2+(m_\theta-m_\tau)\frac{\sqrt\Delta}2=0$,
which gives the formula.
\end{proof}

A \emph{Moore graph} of diameter $2$ is a $k$-regular graph of diameter $2$ with the largest possible number of vertices: from a
vertex, there are $k$ neighbours and at most $k(k-1)$ vertices at distance $2$, so $n\le k^2+1$, and equality means that there are no
triangles and no $4$-cycles (Exercise~\ref{linear:ex:moorebound}). Then two adjacent vertices have no common neighbour, and two non-adjacent
ones have exactly one: a Moore graph is strongly regular with parameters $(k^2+1,k,0,1)$.

\begin{theorem}[Hoffman--Singleton, 1960]\label{linear:thm:moore}
A Moore graph of diameter $2$ and degree $k$ exists only if $k\in\{2,3,7,57\}$.
\end{theorem}
\begin{proof}
Here $\lambda-\mu=-1$, $k-\mu=k-1$, $\Delta=4k-3$, and the equations of the proof of Theorem~\ref{linear:thm:srg} read $m_\theta+m_\tau=k^2$
and
\[
(m_\theta-m_\tau)\sqrt{4k-3}=k^2-2k.
\]
If $\sqrt{4k-3}$ is irrational, then $m_\theta=m_\tau$ and $k^2=2k$, so $k=2$: the pentagon. Otherwise $4k-3=s^2$ with $s$ a positive
odd integer, $k=\frac{s^2+3}4$, and substituting, $16(m_\theta-m_\tau)s=(s^2+3)(s^2-5)=s^4-2s^2-15$. So $s$ divides $15$:
$s\in\{1,3,5,15\}$ and $k\in\{1,3,7,57\}$. For $k=1$ the graph is a single edge, of diameter $1$, not $2$.
\end{proof}

The case $k=3$ is the Petersen graph, and $k=7$ is the \emph{Hoffman--Singleton graph}, on $50$ vertices, unique, which a numerical calculation
builds from five pentagons and five pentagrams. Whether a Moore graph of degree $57$ exists, with $3250$ vertices, is one of the
famous open problems of algebraic combinatorics; Aschbacher (1971) proved it could not be vertex-transitive, and much more is now
known about its possible symmetries, but neither a construction nor a proof of impossibility is in sight. It reappears in
Tome~15.

\section{Independent sets and colourings}

A set of vertices is \emph{independent} if it contains no edge; $\alpha(G)$ is the largest size of an independent set. Computing
$\alpha$ is NP-hard in general, so bounds computable in polynomial time are precious.

\begin{theorem}[Hoffman's ratio bound]\label{linear:thm:hoffman}
Let $G$ be $k$-regular on $n$ vertices with smallest eigenvalue $\lambda_n<0$. Then
\[
\alpha(G)\le n\,\frac{-\lambda_n}{k-\lambda_n}.
\]
\end{theorem}
\begin{proof}
Let $S$ be independent, $|S|=s$, and $x=\mathbf 1_S-\frac sn\mathbf 1$, orthogonal to $\mathbf 1$. Since $\mathbf 1_S^{\mathsf T}A\mathbf 1_S=2\,
|E(S)|=0$ and $A\mathbf 1=k\mathbf 1$,
\[
x^{\mathsf T}Ax=0-2\frac sn\,\mathbf 1_S^{\mathsf T}A\mathbf 1+\frac{s^2}{n^2}\mathbf 1^{\mathsf T}A\mathbf 1=-2\frac{s^2k}n+\frac{s^2k}n=-\frac{s^2k}n,
\]
and $\|x\|^2=s-\frac{s^2}n=\frac{s(n-s)}n$. The Rayleigh quotient is at least $\lambda_n$, so $-\frac{s^2k}n\ge\lambda_n\frac{s(n-s)}n$,
that is $s(k-\lambda_n)\le-\lambda_nn$.
\end{proof}

For the Petersen graph, the bound is $10\cdot\frac2{3+2}=4$, attained by the four $2$-subsets of $\{1,\dots,5\}$ containing $1$.
That example generalises.

\begin{example}[Erd\H{o}s--Ko--Rado for pairs]\label{linear:exa:ekr}
The \emph{Kneser graph} $K(n,2)$ has the $\binom n2$ pairs from $\{1,\dots,n\}$ as vertices, two pairs adjacent when disjoint. An
independent set is a family of pairwise intersecting pairs. Its spectrum can be computed from the incidence matrix $B$ ($n\times
\binom n2$, $b_{i,e}=1$ if $i\in e$): $BB^{\mathsf T}=(n-2)I+J$ has eigenvalues $2(n-1)$ and $n-2$ ($n-1$ times), so $B^{\mathsf T}B$
has these and $0$ for the remaining $\binom n2-n$ dimensions. Two distinct pairs meet in exactly one point if and only if
$(B^{\mathsf T}B)_{ef}=1$, so the ``meeting'' graph has adjacency $B^{\mathsf T}B-2I$, with eigenvalues $2n-4$, $n-4$, $-2$; the
Kneser graph is its complement, and on $\mathbf 1^\perp$ its adjacency $J-I-(B^{\mathsf T}B-2I)$ has eigenvalues $-1-(n-4)=3-n$ and
$-1+2=1$, while its degree is $\binom{n-2}2$. For $n\ge4$, Hoffman's bound gives
\[
\alpha\le\binom n2\frac{n-3}{\binom{n-2}2+n-3}=\binom n2\frac{n-3}{\frac{n(n-3)}2}=n-1.
\]
So a family of pairwise intersecting pairs has at most $n-1$ members: the case $k=2$ of the Erd\H{o}s--Ko--Rado theorem, whose
general case is a classical theorem of extremal set theory.
\end{example}

The same idea bounds the chromatic number: a proper colouring splits $V$ into $\chi(G)$ independent sets, so
$\chi(G)\ge n/\alpha(G)\ge1+\frac k{-\lambda_n}$ for regular $G$ (Exercise~\ref{linear:ex:chromatic} removes regularity).

\section{The Graham--Pollak theorem}

In the 1970s, Ron Graham and Henry Pollak at Bell Labs were designing addressing schemes for telephone networks and needed to
cover the edges of a complete graph by complete bipartite graphs (``bicliques'') with no edge covered twice. The $n-1$ stars at
the vertices $1,\dots,n-1$ (star $i$ joins $i$ to all $j>i$) do it. Can one do better?

\begin{theorem}[Graham--Pollak, 1971]\label{linear:thm:grahampollak}
The edge set of $K_n$ cannot be partitioned into fewer than $n-1$ complete bipartite graphs.
\end{theorem}
\begin{proof}
(Tverberg's proof.) Suppose the edges are partitioned into bicliques with parts $(X_1,Y_1),\dots,(X_m,Y_m)$. Every edge $\{i,j\}$
lies in exactly one of them, with one end in $X_t$ and the other in $Y_t$, so for every $x\in\R^n$
\[
\sum_{i<j}x_ix_j=\sum_{t=1}^m\Bigl(\sum_{i\in X_t}x_i\Bigr)\Bigl(\sum_{j\in Y_t}x_j\Bigr).
\]
If $m<n-1$, the $m+1<n$ homogeneous linear equations $\sum_{i\in X_t}x_i=0$ ($t=1,\dots,m$) and $\sum_ix_i=0$ have a non-zero solution
$x$. For it, the right side vanishes, so $\sum_{i<j}x_ix_j=0$, and $0=\bigl(\sum_ix_i\bigr)^2=\sum_ix_i^2+2\sum_{i<j}x_ix_j=\sum_i
x_i^2>0$, a contradiction.
\end{proof}

The theorem is also an eigenvalue statement in disguise: $J-I$ has $n-1$ negative eigenvalues, and each biclique contributes a
matrix with exactly one (Exercise~\ref{linear:ex:inertia}). No purely combinatorial proof is known: every proof found so far uses linear
algebra.

\section{Huang's theorem}

The hypercube $Q_n$ has $2^n$ vertices and is bipartite, the two sides being the words of even and odd weight, each of size
$2^{n-1}$. So an induced subgraph on $2^{n-1}$ vertices can have no edges at all. Add one more vertex, and the question becomes:
how small can the maximum degree be? In 1988, Chung, F\"uredi, Graham and Seymour constructed examples with maximum degree about
$\sqrt n$, and in 1992 Gotsman and Linial showed that a lower bound of the form $n^c$, for any $c>0$, would prove the sensitivity
conjecture of Nisan and Szegedy (see the remark below).

\begin{theorem}[Huang, 2019]\label{linear:thm:huang}
Every induced subgraph of $Q_n$ on $2^{n-1}+1$ vertices has a vertex of degree at least $\sqrt n$.
\end{theorem}
\begin{proof}
Define symmetric matrices by $A_1=\bigl(\begin{smallmatrix}0&1\\1&0\end{smallmatrix}\bigr)$ and
\[
A_n=\begin{pmatrix}A_{n-1}&I\\I&-A_{n-1}\end{pmatrix}\in M_{2^n}(\Z).
\]
\emph{(i) $A_n^2=nI$.} By induction, since
\[
A_n^2=\begin{pmatrix}A_{n-1}^2+I&A_{n-1}-A_{n-1}\\A_{n-1}-A_{n-1}&I+A_{n-1}^2\end{pmatrix}=nI.
\] So the eigenvalues of $A_n$ are $\pm\sqrt n$, and since $\operatorname{tr}A_n=\operatorname{tr}A_{n-1}-
\operatorname{tr}A_{n-1}=0$, each has multiplicity $2^{n-1}$.

\emph{(ii) $A_n$ is a signed adjacency matrix of $Q_n$}: its non-zero entries are $\pm1$, exactly at the positions of the edges of $Q_n$,
when the vertices are ordered so that the first half has last coordinate $0$. Indeed $Q_n$ consists of two copies of $Q_{n-1}$ (the two
diagonal blocks) joined by a perfect matching between corresponding vertices (the identity blocks).

\emph{(iii) Interlacing.} Let $H$ be induced on a set $S$ of $2^{n-1}+1$ vertices, and $B$ the principal submatrix of $A_n$ on $S$. By
Cauchy's interlacing theorem (Theorem~\ref{linear:thm:interlacing} with $m=2^{n-1}+1$),
\[
\lambda_1(B)\ge\lambda_{1+2^n-m}(A_n)=\lambda_{2^{n-1}}(A_n)=\sqrt n.
\]
\emph{(iv) Degree bound.} Let $By=\lambda_1(B)y$ and $|y_i|$ be the largest coordinate. Then $\lambda_1(B)|y_i|=\bigl|\sum_jb_{ij}y_j\bigr|\le
\sum_j|b_{ij}||y_i|=\deg_H(i)|y_i|$, so $\Delta(H)\ge\lambda_1(B)\ge\sqrt n$.
\end{proof}

\begin{figure}[ht]
\centering
\fig{linear-ch08-huang}
\caption{Huang's signing of the cube $Q_3$: dashed red edges carry $-1$. Every square face has an odd number of negative edges,
which is what makes the signed matrix square to $3I$.}
\label{linear:fig:huang}
\end{figure}

\begin{remark}
A Boolean function $f:\{0,1\}^n\to\{0,1\}$ has \emph{sensitivity} $s(f)$, the largest number of coordinates whose flip changes
$f(x)$, over all inputs $x$; and \emph{block sensitivity} $bs(f)$, the same with disjoint blocks of coordinates flipped together.
Always $s(f)\le bs(f)$, and Nisan and Szegedy asked in 1992 whether $bs(f)\le s(f)^C$ for a constant $C$. Gotsman and Linial had
shown that this would follow from a statement like Theorem~\ref{linear:thm:huang}, and Huang's theorem gives $bs(f)\le2s(f)^4$. Since block
sensitivity is polynomially related to many other complexity measures (decision tree depth, certificate complexity, degree as a
polynomial), sensitivity joined them. The reduction is combinatorial and is told in Huang's paper and in Aaronson's and Tao's blog
posts listed in the sources.
\end{remark}

\section{A gem: three Petersens do not make a $K_{10}$}

The complete graph $K_{10}$ has $45$ edges, and the Petersen graph $15$. Can the edges of $K_{10}$ be split into three copies of the
Petersen graph? Counting says maybe; eigenvalues say no.

\begin{proposition}[Schwenk]\label{linear:prop:schwenk}
$K_{10}$ is not the edge-disjoint union of three Petersen graphs.
\end{proposition}
\begin{proof}
Suppose $J-I=P_1+P_2+P_3$ with each $P_t$ the adjacency matrix of a Petersen graph (on the same ten vertices, differently numbered).
Each $P_t$ is $3$-regular, and its eigenspace for the eigenvalue $1$ is a $5$-dimensional subspace of the $9$-dimensional space $\mathbf
1^\perp$. Two $5$-dimensional subspaces of a $9$-dimensional space intersect non-trivially, so there is $x\ne0$, $x\perp\mathbf 1$,
with $P_1x=P_2x=x$. Then $P_3x=(J-I)x-P_1x-P_2x=-x-x-x=-3x$, so $-3$ would be an eigenvalue of the Petersen graph $P_3$. It is not.
\end{proof}

\begin{beauty}
Graphs are combinatorial objects, and yet some of their deepest properties are forced by arithmetic and geometry of real numbers:
a square root that cannot be rational (friendship), an integer that must divide $15$ (Moore graphs), two subspaces too large to
miss each other (Petersen), a corner of a matrix that cannot escape interlacing (Huang). The spectrum is not a complete
fingerprint (different graphs can share one), but when the fingerprint is impossible, so is the graph.
\end{beauty}

\section{Exercises}

\begin{exercise}[$\star$]\label{linear:ex:hypercube}
Show that the functions $\chi_T(x)=(-1)^{\sum_{i\in T}x_i}$, $T\subseteq\{1,\dots,n\}$, are eigenvectors of the adjacency matrix of
$Q_n$ with eigenvalues $n-2|T|$.
\end{exercise}

\begin{exercise}[$\star$]\label{linear:ex:spectrumcount}
A graph has spectrum $3,1,1,1,1,1,-2,-2,-2,-2$. How many vertices, edges and triangles does it have? Is it regular?
\end{exercise}

\begin{exercise}[$\star\star$]\label{linear:ex:bipartite}
Let $G$ be connected. Show that $G$ is bipartite if and only if $-\lambda_1$ is an eigenvalue, if and only if the spectrum is
symmetric about $0$.
\end{exercise}

\begin{exercise}[$\star$]\label{linear:ex:srgcount}
Show that the parameters of a strongly regular graph satisfy $k(k-\lambda-1)=(n-k-1)\mu$, and compute the eigenvalues and
multiplicities of the Petersen graph and of the hypothetical Moore graph of degree $57$.
\end{exercise}

\begin{exercise}[$\star$]\label{linear:ex:moorebound}
Show that a graph of diameter $2$ and maximum degree $k$ has at most $k^2+1$ vertices, with equality only if it is $k$-regular
without triangles and $4$-cycles.
\end{exercise}

\begin{exercise}[$\star\star$]\label{linear:ex:linegraph}
The \emph{line graph} $L(G)$ has the edges of $G$ as vertices, adjacent when they share an endpoint. Show that its adjacency matrix
is $B^{\mathsf T}B-2I$, with $B$ the vertex-edge incidence matrix of $G$, and deduce that every eigenvalue of a line graph is at least
$-2$.
\end{exercise}

\begin{exercise}[$\star\star$]\label{linear:ex:chromatic}
(Hoffman.) Show that every graph with at least one edge satisfies $\chi(G)\ge1+\frac{\lambda_1}{-\lambda_n}$. (Test the Rayleigh quotient on the
Perron vector, twisted by roots of unity along the colour classes.)
\end{exercise}

\begin{exercise}[$\star\star$]\label{linear:ex:wilf}
(Wilf.) Show that $\chi(G)\le1+\lambda_1$ for every graph. (Use that some vertex has degree at most the average degree, which is at
most $\lambda_1$, and interlacing.)
\end{exercise}

\begin{exercise}[$\star\star\star$]\label{linear:ex:inertia}
Show that a biclique on parts $X,Y$ has adjacency matrix $\mathbf 1_X\mathbf 1_Y^{\mathsf T}+\mathbf 1_Y\mathbf 1_X^{\mathsf T}$, with exactly
one negative eigenvalue, and that the number of negative eigenvalues is subadditive: $\nu(A+B)\le\nu(A)+\nu(B)$ for Hermitian $A,B$.
Deduce Theorem~\ref{linear:thm:grahampollak} again.
\end{exercise}

\begin{exercise}[$\star\star$]\label{linear:ex:infinitefriends}
Build an infinite graph in which every two vertices have exactly one common neighbour and no vertex is adjacent to all others.
(Start from three vertices and no edges, and repeatedly add, for each pair without a common neighbour, a new vertex adjacent to both.)
\end{exercise}

\begin{exercise}[$\star\star$]\label{linear:ex:huangsharp}
Show that Theorem~\ref{linear:thm:huang} is sharp for $n=1$ and $n=4$: find an induced subgraph of $Q_4$ on $9$ vertices with maximum
degree $2$. (Hint: start from the even-weight words and replace a few of them.)
\end{exercise}

\section{Hints}

\begin{hint}{linear:ex:hypercube}
$(A\chi_T)(x)=\sum_{i=1}^n\chi_T(x+e_i)$, and flipping coordinate $i$ multiplies $\chi_T$ by $-1$ if $i\in T$, by $1$ otherwise.
\end{hint}
\begin{hint}{linear:ex:spectrumcount}
Use $\operatorname{tr}A^0$, $\operatorname{tr}A^2$, $\operatorname{tr}A^3$; for regularity, compare $\lambda_1$ with the average degree.
\end{hint}
\begin{hint}{linear:ex:bipartite}
For a bipartite graph, $D=\operatorname{diag}(\pm1)$ by side satisfies $DAD=-A$. Conversely, look at the proof of Frobenius's
theorem for the eigenvalue $-\lambda_1$.
\end{hint}
\begin{hint}{linear:ex:srgcount}
Count the pairs $(x,y)$ with $u\sim x\sim y$, $u\not\sim y$, $y\ne u$, in two ways.
\end{hint}
\begin{hint}{linear:ex:moorebound}
Count vertices at distance $1$ and $2$ from a fixed vertex.
\end{hint}
\begin{hint}{linear:ex:linegraph}
$(B^{\mathsf T}B)_{ef}$ is the number of common endpoints of the edges $e,f$.
\end{hint}
\begin{hint}{linear:ex:chromatic}
If $V=V_1\sqcup\dots\sqcup V_c$ with independent $V_t$, cut a unit Perron vector $x$ into its pieces $x_t$ on the classes $V_t$.
Twist the pieces by roots of unity, $y_j=\sum_t\omega^{jt}x_t$ with $\omega=e^{2\pi i/c}$, compare each $y_j^*Ay_j$ with $\lambda_n$,
and average over $j=1,\dots,c-1$.
\end{hint}
\begin{hint}{linear:ex:wilf}
Induction on $n$: remove a vertex of minimum degree.
\end{hint}
\begin{hint}{linear:ex:inertia}
Weyl's inequality in the form $\lambda_{i+j-1}(-A-B)\le\lambda_i(-A)+\lambda_j(-B)$.
\end{hint}
\begin{hint}{linear:ex:infinitefriends}
Make sure the new vertex creates no second common neighbour for any pair.
\end{hint}
\begin{hint}{linear:ex:huangsharp}
Try a finite search for one.
\end{hint}

\section{Solutions}

\begin{solution}{linear:ex:hypercube}
\idea Characters of $(\Z/2)^n$ diagonalise every ``translation-invariant'' operator; the adjacency operator of $Q_n$ is a sum of
$n$ translations.
\writeup The neighbours of $x$ are $x+e_i$ (addition mod $2$). Since $\chi_T(x+e_i)=-\chi_T(x)$ if $i\in T$ and $\chi_T(x)$ otherwise,
$(A\chi_T)(x)=\sum_i\chi_T(x+e_i)=(n-|T|-|T|)\chi_T(x)$. The $2^n$ functions $\chi_T$ are orthogonal ($\sum_x\chi_T(x)\chi_{T'}(x)=
\sum_x\chi_{T\triangle T'}(x)=0$ unless $T=T'$), so they form an eigenbasis, and $n-2j$ has multiplicity $\binom nj$.
\end{solution}

\begin{solution}{linear:ex:spectrumcount}
\idea Traces of powers count closed walks.
\writeup $n=10$; $\operatorname{tr}A^2=9+5+16=30=2|E|$, so $15$ edges; $\operatorname{tr}A^3=27+5-32=0$, so no triangles. The
average degree is $3=\lambda_1$, so the graph is $3$-regular (equality case of Exercise~\ref{linear:ex:degree}). It is the spectrum of the
Petersen graph (and indeed only the Petersen graph has it, though proving that takes more work).
\end{solution}

\begin{solution}{linear:ex:bipartite}
\idea Bipartite graphs have period $2$ in the sense of Chapter~\ref{linear:ch:perron}; the signs of an eigenvector for $-\lambda_1$ reveal the two
sides.
\writeup If $G$ is bipartite with sides $V_1,V_2$, $D=\operatorname{diag}(d_v)$ with $d_v=1$ on $V_1$ and $-1$ on $V_2$ satisfies
$DAD=-A$, so $A$ and $-A$ are similar: the spectrum is symmetric, and in particular contains $-\lambda_1$. Conversely, let $Ay=-\lambda_1y$.
The proof of Theorem~\ref{linear:thm:frobenius} (with $\varphi=\pi$) shows $|y|>0$ and $e^{i\theta_j}=-e^{i\theta_i}$ whenever $i\sim j$; so the
sets $\{v:e^{i\theta_v}=e^{i\theta_{v_0}}\}$ and $\{v:e^{i\theta_v}=-e^{i\theta_{v_0}}\}$ cover $V$ (by connectivity) and every edge joins
them: $G$ is bipartite. The third condition implies the second, and the first implies the third.
\end{solution}

\begin{solution}{linear:ex:srgcount}
\idea Double counting of paths of length $2$ leaving the neighbourhood.
\writeup Fix $u$. Each of the $k$ neighbours $x$ of $u$ has $k-1-\lambda$ neighbours other than $u$ that are not neighbours of $u$; each
of the $n-k-1$ non-neighbours $y$ of $u$ has $\mu$ neighbours among the neighbours of $u$. So $k(k-1-\lambda)=(n-k-1)\mu$. Petersen
$(10,3,0,1)$: $x^2+x-2=0$, $\theta=1$, $\tau=-2$, $\Delta=9$, multiplicities $\frac12(9\mp\frac{6-9}3)=5,4$. Moore graph of degree $57$:
$n=3250$, $\Delta=225$, $\theta=7$, $\tau=-8$, and $m_\theta,m_\tau=\frac12\bigl(3249\mp\frac{114-3249}{15}\bigr)=\frac12(3249\pm209)$, that is
$1729$ and $1520$. (Yes, $1729$, Ramanujan's taxicab number.)
\end{solution}

\begin{solution}{linear:ex:moorebound}
\idea Breadth-first search from one vertex: the layers have at most $1$, $k$ and $k(k-1)$ vertices.
\writeup Fix $u$. It has at most $k$ neighbours, each with at most $k-1$ further neighbours, and by diameter $2$ every vertex is $u$, a
neighbour, or a neighbour of a neighbour: $n\le1+k+k(k-1)=k^2+1$. Equality forces $\deg u=k$, every neighbour of degree $k$, and the
$k(k-1)$ second neighbours all distinct and different from $u$ and its neighbours. A triangle $u,x,x'$ would make $x'$ counted as a
neighbour of $x$ though already a neighbour of $u$; a $4$-cycle $u,x,y,x'$ would count $y$ twice. As $u$ is arbitrary, all degrees are
$k$ and there are no triangles or $4$-cycles.
\end{solution}

\begin{solution}{linear:ex:linegraph}
\idea Incidence matrices turn ``share an endpoint'' into an inner product of columns.
\writeup Column $e$ of $B$ is $\mathbf 1_{\{i,j\}}$ for $e=\{i,j\}$, so $(B^{\mathsf T}B)_{ef}=|e\cap f|$, which is $2$ on the diagonal, $1$
for distinct edges sharing an endpoint and $0$ otherwise. Hence $A(L(G))=B^{\mathsf T}B-2I$, and since $B^{\mathsf T}B\succeq0$, all its
eigenvalues are $\ge-2$.
\end{solution}

\begin{solution}{linear:ex:chromatic}
\idea Use the colour classes to cut the matrix into blocks with zero diagonal blocks, and let interlacing-type inequalities compare
$\lambda_1$ with $\lambda_n$.
\writeup Let $V=V_1\sqcup\dots\sqcup V_c$ be a proper colouring and $x$ a unit Perron vector of $A$ ($x\ge0$, $Ax=\lambda_1x$). Let
$x_t=x\cdot\mathbf 1_{V_t}$ (zero outside $V_t$) and $s_t=\|x_t\|^2$, so $\sum_ts_t=1$. Since $V_t$ is independent, $x_t^{\mathsf T}Ax_t=0$.
For any signs, consider $y=\sum_t\varepsilon_tx_t$ with $\varepsilon_t\in\C$, $|\varepsilon_t|=1$: then $y^*Ay=\sum_{t\ne t'}\bar\varepsilon_t
\varepsilon_{t'}x_t^{\mathsf T}Ax_{t'}$, with $\|y\|=1$, so $y^*Ay\ge\lambda_n$. Take $\varepsilon_t=\omega^{t}$ with $\omega=e^{2\pi ij/c}$ and
average over $j=1,\dots,c-1$: since $\sum_{j=0}^{c-1}\omega^{j(t'-t)}=0$ for $t\ne t'$, the average over $j=1,\dots,c-1$ of
$\bar\varepsilon_t\varepsilon_{t'}$ is $-\frac1{c-1}$. Hence
\[
\lambda_n\le-\frac1{c-1}\sum_{t\ne t'}x_t^{\mathsf T}Ax_{t'}=-\frac1{c-1}x^{\mathsf T}Ax=-\frac{\lambda_1}{c-1},
\]
using $x^{\mathsf T}Ax=\sum_{t,t'}x_t^{\mathsf T}Ax_{t'}$ and $x_t^{\mathsf T}Ax_t=0$. So $c-1\ge\lambda_1/(-\lambda_n)$.
\end{solution}

\begin{solution}{linear:ex:wilf}
\idea Greedy colouring needs at most ``degree $+1$'' colours; the right degree to use is not the maximum but $\lambda_1$, thanks to
interlacing.
\writeup By induction on $n$. Some vertex $v$ has $\deg v\le\frac{2|E|}n=R_A(\mathbf 1)\le\lambda_1(G)$. The graph $G-v$ has adjacency
matrix a principal submatrix of $A$, so $\lambda_1(G-v)\le\lambda_1(G)$ by interlacing; by induction it can be coloured with
$1+\lambda_1(G-v)\le1+\lambda_1(G)$ colours (taking integer parts: $\lfloor1+\lambda_1\rfloor$ colours). Then $v$ has at most
$\lfloor\lambda_1(G)\rfloor$ neighbours, so one of $\lfloor1+\lambda_1(G)\rfloor$ colours is free for $v$.
\end{solution}

\begin{solution}{linear:ex:inertia}
\idea A biclique is a rank-two matrix $uv^{\mathsf T}+vu^{\mathsf T}$ with $u\perp v$, which has one positive and one negative eigenvalue;
negative eigenvalues cannot be created faster than one per summand.
\writeup With $u=\mathbf 1_X$, $v=\mathbf 1_Y$ ($X\cap Y=\emptyset$, so $u\perp v$), $M=uv^{\mathsf T}+vu^{\mathsf T}$ maps $\frac u{\|u\|}\pm
\frac v{\|v\|}$ to $\pm\|u\|\|v\|(\frac u{\|u\|}\pm\frac v{\|v\|})$ and kills $\{u,v\}^\perp$: exactly one negative eigenvalue. For
subadditivity, $\nu(M)$ is the number of positive eigenvalues of $-M$. If $\nu(A)=a$ and $\nu(B)=b$, then $\lambda_{a+1}(-A)\le0$ and
$\lambda_{b+1}(-B)\le0$, and Weyl gives $\lambda_{a+b+1}(-A-B)\le\lambda_{a+1}(-A)+\lambda_{b+1}(-B)\le0$: at most $a+b$ eigenvalues of
$-(A+B)$ are positive. If $J-I=\sum_{t=1}^mM_t$, then $n-1=\nu(J-I)\le\sum_t\nu(M_t)=m$.
\end{solution}

\begin{solution}{linear:ex:infinitefriends}
\idea Fix every pair without a common neighbour by a new private vertex; this never creates a second common neighbour, and the new
pairs are fixed at the next stage. The limit is a friendship graph, and degrees grow too fast for a politician.
\writeup Start with $G_0$: three vertices $a,b,c$ and no edges. Given $G_k$, for each pair $\{u,v\}$ of vertices of $G_k$ with no
common neighbour in $G_k$, add a new vertex $w_{uv}$ adjacent to $u$ and $v$ only; call $\{u,v\}$ its parent pair. Let us check
that in $G_{k+1}$ every pair has at most one common neighbour, assuming this holds in $G_k$.
\begin{itemize}
\item Two old vertices: a pair that had a common neighbour gets no new one, since every new vertex is attached to a pair that had
none; a pair that had none gets exactly one, its own $w_{uv}$.
\item A new vertex $w_{uv}$ and an old vertex $x$: a common neighbour lies in $\{u,v\}$, and it cannot be both $u$ and $v$, because
then $x$ would have been a common neighbour of $u$ and $v$.
\item Two new vertices: a common neighbour lies in the intersection of the two parent pairs, which has at most one element.
\end{itemize}
So in $G=\bigcup_kG_k$ every pair has at most one common neighbour (a common neighbour of two vertices of $G_k$ found in $G$ is
already present in some $G_m$), and every pair of vertices of $G_k$ has one in $G_{k+1}$: exactly one.

No vertex is adjacent to all others. Step~4 of the proof of Theorem~\ref{linear:thm:friendship} uses no finiteness: if some $w$ were
adjacent to all others, every other vertex would have degree exactly $2$. But $G_1$ is the hexagon $a,w_{ab},b,w_{bc},c,w_{ca}$; in
it the two edges at $a$ and the opposite pair $\{a,w_{bc}\}$ have no common neighbour, so $a$ gets three new neighbours and has
degree $5$ in $G_2$; likewise $b$. Whichever vertex $w$ were, one of $a,b$ would be a vertex other than $w$ with
degree above $2$. (So which step of the proof breaks for infinite graphs? Step~3, the only one that takes a trace.)
\end{solution}

\begin{solution}{linear:ex:huangsharp}
\idea Nine vertices with every degree at most $2$ means a disjoint union of paths and cycles inside the cube. An $8$-cycle in $Q_4$
plus one far-away vertex is the natural thing to try.
\writeup For $n=1$: $Q_1$ is an edge, and its $2^0+1=2$ vertices induce a graph of maximum degree $1=\sqrt1$. For $n=4$, take the
$8$-cycle
\[
0000-0001-0101-0111-1111-1011-1010-0010-0000
\]
together with the vertex $1100$. Each step of the cycle flips one bit, and no other two of the eight words differ in exactly one
bit (check the $28$ pairs). The neighbours of $1100$ are $0100,1000,1101,1110$, none of which is in the set. So the induced
subgraph has maximum degree exactly $2=\sqrt4$, and Huang's bound is attained. Numerical evaluation confirms it by brute force: exactly
$48$ of the $\binom{16}9=11440$ nine-element sets have maximum degree $2$, and every one of them is an $8$-cycle plus an isolated
vertex.
\end{solution}

\section*{Sources and further reading}

\begin{itemize}
\item P.~Erd\H{o}s, A.~R\'enyi and V.~T.~S\'os, \emph{On a problem of graph theory}, Studia Scientiarum Mathematicarum Hungarica 1
(1966), 215--235; M.~Aigner and G.~M.~Ziegler, \emph{Proofs from THE BOOK}, 6th ed., Springer, 2018, chapters ``Of friends and
politicians'' and ``Three famous theorems on finite sets'' (Graham--Pollak).
\item A.~J.~Hoffman and R.~R.~Singleton, \emph{On Moore graphs with diameters 2 and 3}, IBM Journal of Research and Development 4 (1960),
497--504; M.~Aschbacher, \emph{The non-existence of rank three permutation groups of degree 3250 and subdegree 57}, Journal of Algebra 19
(1971), 538--540.
\item A.~J.~Hoffman, \emph{On eigenvalues and colorings of graphs}, in \emph{Graph Theory and its Applications}, Academic Press, 1970,
79--91; H.~S.~Wilf, \emph{The eigenvalues of a graph and its chromatic number}, Journal of the London Mathematical Society 42 (1967),
330--332.
\item R.~L.~Graham and H.~O.~Pollak, \emph{On the addressing problem for loop switching}, Bell System Technical Journal 50 (1971),
2495--2519; H.~Tverberg, \emph{On the decomposition of $K_n$ into complete bipartite graphs}, Journal of Graph Theory 6 (1982), 493--494.
\item A.~J.~Schwenk, \emph{Problem 6434}, American Mathematical Monthly 90 (1983), 403, and solution 92 (1985), 591--592 (three Petersens).
\item H.~Huang, \emph{Induced subgraphs of hypercubes and a proof of the Sensitivity Conjecture}, Annals of Mathematics 190 (2019),
949--955; F.~R.~K.~Chung, Z.~F\"uredi, R.~L.~Graham and P.~Seymour, \emph{On induced subgraphs of the cube}, Journal of Combinatorial
Theory A 49 (1988), 180--187; C.~Gotsman and N.~Linial, \emph{The equivalence of two problems on the cube}, Journal of Combinatorial Theory
A 61 (1992), 142--146; N.~Nisan and M.~Szegedy, \emph{On the degree of Boolean functions as real polynomials}, Computational Complexity 4
(1994), 301--313.
\item S.~Aaronson, \emph{Sensitivity conjecture resolved}, blog post, 2 July 2019; T.~Tao, \emph{Twisted convolution and the sensitivity
conjecture}, blog post, 2 July 2019.
\item A.~E.~Brouwer and W.~H.~Haemers, \emph{Spectra of Graphs}, Springer, 2012; C.~Godsil and G.~Royle, \emph{Algebraic Graph Theory},
Springer, 2001, Chapters 5, 9 and 10.
\end{itemize}
